\begin{textsample}{NEG Dim 1 – persona\_gemini – Score -1.00 – t475\_gemini.txt}  \label{ex:f1_neg_013}
A Greenpeace East Asia report has exposed links between Taiwanese seafood giant Fong Chun Formosa ( FCF ) and severe abuses in the fishing industry.
FCF, a top tuna trader with an annual revenue of \$45 billion, has been connected to forced labor on migrant workers, illegal fishing, \textbf{shark} finning, and transshipment at sea.
The exploitation is enabled by Taiwan 's two-tiered system, which excludes vulnerable Southeast Asian \textbf{fishers} from proper labor protections.
We urge the Taiwanese government to abolish this system, adopt International Labour Organization ( ILO ) conventions, enhance inspections, and create grievance mechanisms for workers.
We also call on FCF to take immediate action by disclosing its \textbf{suppliers}, ending the practice of transshipment, ensuring adequate \textbf{monitoring}, and upgrading its human rights policies.
Consumers can help combat these industry abuses by reducing their seafood intake and supporting local fisheries.

% matched lemmas: fisher, monitoring, shark, supplier
\end{textsample}
